% The usual "AB" scheme: two sets of 19 arrays, A (f, the values f_i) and B (g, the new values f'_i), drawn as the
% two field rows of the slide "How the 19 values are stored in memory" (figures/fzyx.tex). Each step reads one set
% and writes the other; then they swap roles (lbm_bulk.py: swap_fields(pdfs, pdfs_tmp) after every sweep).
% presentation7-v3: A and B look the same as where they are introduced.
\begin{tikzpicture}[x=1cm,y=1cm]
  \tikzset{ca/.style={draw=psgray!70,line width=0.3pt,minimum height=4.2mm,minimum width=2.7mm,inner sep=0pt,
             font=\tiny,anchor=west},
           cb/.style={ca,fill=psorange!12},
           st/.style={-{Stealth[length=1.8mm]},thick}}
  \def\w{0.27}
  %   \fieldrow{y}{style}{prime}: 19 arrays, the first three and the last labelled
  \newcommand{\fieldrow}[3]{%
    \foreach \i in {0,...,18} {
      \ifnum\i<2 \node[#2] at (\i*\w,#1) {$f#3_{\i}$};
      \else\ifnum\i=18 \node[#2] at (\i*\w,#1) {\tiny$f#3_{18}$};
      \else \node[#2] at (\i*\w,#1) {};\fi\fi}}
  \fieldrow{0}{ca}{}
  \fieldrow{-2.0}{cb}{'}
  \node[font=\small\bfseries,anchor=east] at (-0.1,0) {A};
  \node[font=\small\bfseries,anchor=east] at (-0.1,-2.0) {B};
  \draw[st,psblue!85!black] (1.3,-0.3) -- node[left,font=\scriptsize,align=right,text=black] {step 1:\\read A,\\write B} (1.3,-1.7);
  \draw[st,psorange!90!black] (3.9,-1.7) -- node[right,font=\scriptsize,align=left,text=black] {step 2:\\read B,\\write A} (3.9,-0.3);
  \node[lbl,anchor=north,align=center] at (19*\w/2,-2.35) {two full copies of the 19 arrays;\\they swap roles after every step};
\end{tikzpicture}
